\section{Related Work}
\label{sec:related}

\noindent\textbf{Pre-trained model-based exemplar-free CIL.}
The use of pre-trained vision models has recently reshaped exemplar-free CIL. Rather than repeatedly updating the full backbone, these methods adapt lightweight components or classifiers on top of transferable representations.
L2P~\cite{wang2022learning} is an early representative work in this direction, which keeps the pre-trained model frozen and learns a prompt pool that can be retrieved without task identities at test time.
DualPrompt~\cite{wang2022dualprompt} further separates prompts into general and expert prompts to improve the stability-plasticity trade-off.
CODA-Prompt~\cite{smith2023coda} extends this line by composing prompt components through an attention-based mechanism.
Following these works, subsequent prompt-based methods improve prompt selection, composition, regularization, or efficiency~\cite{tang2023prompt,khan2023introducing,kim2024one,kurniawan2024evolving,qiao2024prompt,huang2024ovor,roy2024convolutional,gao2024consistent,lu2024visual,zhang2025vpt}.
These methods show that pre-trained models can support continual learning with only lightweight task-adaptive parameters. However, their performance still depends on whether the appropriate prompts can be retrieved or composed for each test sample.

Another line of work avoids test-time prompt or component selection by using fixed representations with lightweight classifiers or class statistics.
ACIL~\cite{acil2022} updates classifiers in closed form without storing historical samples.
FeCAM~\cite{fecam2023} performs classification with class-wise statistics in the feature space, while RanPAC~\cite{ranpac2023} improves prototype-based classification with random projection and prototype accumulation.
Recent studies further revisit prototype- and classifier-based exemplar-free CIL under stronger pre-trained representations~\cite{zhou2025revisiting,wang2024min}.
These methods are simple and stable because all seen classes are compared under a unified task-agnostic decision rule.
However, the same global rule may provide limited task-specific adaptation when classes from different tasks require different local decision boundaries.

Beyond prompts and shared classifiers, recent works introduce adapters, subspaces, or experts on top of pre-trained models.
Some methods isolate task-specific parameters to reduce interference among incremental tasks~\cite{gao2023unified,wang2023isolation}.
EASE~\cite{ease2024} expands task-specific subspaces with lightweight adapters, while InfLoRA~\cite{liang2024inflora} adopts low-rank adaptation to improve parameter efficiency and stability.
More recent adapter- and expert-based methods further explore component fusion, retrieval, or aggregation during inference~\cite{gao2024beyond,cofima2024,wang2025integrating,mos2025}.
Compared with global classifiers, these methods provide stronger task-specific discrimination, but they also raise a new question: how to reliably use task-adaptive components when task identities are unavailable.

\noindent\textbf{Task-agnostic inference.}
Task-agnostic inference has therefore become a central issue in pre-trained model-based CIL.
Prompt-based methods such as L2P~\cite{wang2022learning}, DualPrompt~\cite{wang2022dualprompt}, and CODA-Prompt~\cite{smith2023coda} address this issue by retrieving or composing prompts according to the input representation.
Adapter- and expert-based methods instead select, fuse, or aggregate task-adaptive modules during inference~\cite{ease2024,liang2024inflora,wang2025integrating,mos2025}.
These strategies improve specialization, but inaccurate routing may activate components learned for mismatched tasks. Directly combining outputs from independently adapted components can be unreliable without proper normalization or calibration.
In contrast, prototype- and classifier-based methods avoid such routing by comparing all classes in a shared feature space~\cite{fecam2023,ranpac2023,zhou2025revisiting}. However, they sacrifice part of the task-specific discrimination offered by adaptive components.

SPIE is designed around this trade-off.
It retains compact task experts for local discrimination, but avoids hard selection of a single expert.
Instead, SPIE estimates soft task weights from a prototype bank built on a frozen ViT backbone and aggregates expert predictions in probability space.
This allows SPIE to benefit from task-specific components while reducing sensitivity to a single incorrect routing decision.
A low-rank distributional prototype replay mechanism further maintains the prototype bank without storing raw data. This enables exemplar-free learning that combines shared-space task estimation with expert-level discrimination.

